% 原书分章 PDF 第 17—24 页：第 16.5—16.6 节与致谢。
\section{实际表现}

我们将 TMLE 和三种 IPW 估计量的结果，与一种粗分析的结果相比较；粗分析旨在比较四种感兴趣强化治疗策略 $d_\theta$ 对应的生存曲线，而不对混杂和选择偏差作任何调整。将基于风险函数的 IPW 估计量的稳定化权重全部设为 1，就能实施这样的分析。表 16.1 给出了用粗分析比较四条生存曲线在 3 年时的推断。该表的推断可与两类结果对照：一类是表 16.2 中采用四种 PS 估计方法实施三种 IPW 估计量的结果；另一类是表 16.1 中基于靶向学习的结果，即对两个干扰参数均用 SL 估计的 TMLE。检查表 16.1 中粗估计量的结果，可发现一些虽然较弱、但与 ACCORD 和 ADVANCE 随机试验结果一致的证据：更积极的治疗启动规则可能具有有益效应。

\begin{table}[htbp]
\centering
\caption{粗估计量和未截断 TMLE 对 12 个季度（$t_0=11$）时六个累积风险差（RD）$\psi^{\theta_1,\theta_2}$ 的推断比较。}
\label{tab:ch16-one}
\begingroup\scriptsize
\setlength{\tabcolsep}{2.0pt}\renewcommand{\arraystretch}{1.2}
\resizebox{\textwidth}{!}{%
\begin{tabular}{@{}ccrrrrrrrrrrr@{}}
\toprule
& & \multicolumn{5}{c}{粗估计} & \multicolumn{6}{c}{TMLE}\\
\cmidrule(lr){3-7}\cmidrule(l){8-13}
$\theta_1$ & $\theta_2$ & $\psi_n^{\theta_1,\theta_2}$ & $\Gamma_n^{\theta_1,\theta_2}$ & $\psi_n^{\theta_1,\theta_2,-}$ & $\psi_n^{\theta_1,\theta_2,+}$ & $p$ & $\psi_n^{\theta_1,\theta_2,*}$ & $\Gamma_n^{\theta_1,\theta_2,*}$ & $\psi_n^{\theta_1,\theta_2,*,-}$ & $\psi_n^{\theta_1,\theta_2,*,+}$ & $p^*$ & RE\\
\midrule
$8.5$ & $8.0$ & $0.0013$ & $0.0015$ & $-0.0016$ & $0.0042$ & $0.386$ & $0.0070$ & $0.0038$ & $-0.0004$ & $0.0143$ & $0.064$ & $1.07$\\
$8.5$ & $7.5$ & $0.0028$ & $0.0030$ & $-0.0032$ & $0.0088$ & $0.357$ & $0.0221$ & $0.0065$ & $0.0093$ & $0.0349$ & $0.001$ & $1.09$\\
$8.5$ & $7.0$ & $0.0335$ & $0.0111$ & $0.0116$ & $0.0553$ & $0.003$ & $0.0386$ & $0.0112$ & $0.0166$ & $0.0606$ & $0.001$ & $1.11$\\
$8.0$ & $7.5$ & $0.0015$ & $0.0027$ & $-0.0037$ & $0.0068$ & $0.574$ & $0.0151$ & $0.0059$ & $0.0035$ & $0.0267$ & $0.011$ & $1.09$\\
$8.0$ & $7.0$ & $0.0322$ & $0.0111$ & $0.0105$ & $0.0538$ & $0.004$ & $0.0316$ & $0.0112$ & $0.0097$ & $0.0535$ & $0.005$ & $1.11$\\
$7.5$ & $7.0$ & $0.0307$ & $0.0108$ & $0.0096$ & $0.0518$ & $0.004$ & $0.0165$ & $0.0103$ & $-0.0038$ & $0.0368$ & $0.110$ & $1.11$\\
\bottomrule
\end{tabular}}
\endgroup
\par\smallskip\begin{minipage}{\textwidth}\footnotesize
TMLE 基于干扰参数的 SL 估计量，即 $g^\theta_{n,t,\mathrm{SL}}$ 和 $Q^\theta_{n,\mathrm{SL}}$ 构造；上标 $*$ 用于将其与粗估计量区别。基于 $g_{n,t,\mathrm{SL}}$ 的风险函数 IPW 估计量（见表 16.2）的标准误与 TMLE 标准误之比，记为 RE。
\end{minipage}
\end{table}


图 16.2 及表 16.2、表 16.1 中的 TMLE 推断，显示了更积极 TI 策略有益效应的强得多的证据，也展示了该估计量在调整时变混杂和选择偏差方面的实际表现。事实上，无论采用何种干扰参数估计方法，TMLE 均表明四条生存曲线在较早时间就已分离，并始终保持一致的顺序，提示更积极的治疗启动规则，即 A1c 阈值逐步降低的规则，其有益效应逐步增大。根据干扰参数 $Q_0^\theta$ 所用估计方法的不同，3 年时的六个累积风险差中，有四个或五个达到统计显著性。

表 16.2 展示了基于风险函数的 IPW 估计量在偏差调整方面的成功表现，同样显示更积极 TI 策略具有有益效应。无论采用何种行动机制估计方法，该 IPW 估计量得到的四条生存曲线均出现分离（结果见 Neugebauer 等，2016），从视觉上看，几乎与图 16.2 中 TMLE 的结果相同。根据行动机制 $g_0^\theta$ 所用估计方法的不同，3 年时六个累积风险差中有三个或四个达到统计显著性。无界 IPW 估计量在偏差调整方面的表现，在很大程度上取决于行动机制的估计方法（结果见 Neugebauer 等，2016）。当与行动机制的 SL 估计结合时，该 IPW 估计量显示的更积极 TI 策略有益效应证据最强。有界 IPW 估计量在偏差调整方面的结果模式（Neugebauer 等，2016）介于无界 IPW 与基于风险函数的 IPW 估计量之间：随着 PS 估计越来越非参数化，曲线按预期方向逐渐更明显、更一致地分离，这类似于无界 IPW 的结果；但四种 PS 估计方法之间的结果差异又不太明显，这类似于基于风险函数的 IPW 的结果。

我们推测，当从无界 IPW 依次转向有界 IPW、再转向基于风险函数的 IPW 估计量时，各 PS 估计方法所得结果模式的逐步变化，与这三种估计量的有界性和有效性差异有关。更具体地，我们推测，这三种估计量对 PS 估计误差所致偏差的敏感性逐步减小，首先归因于有界 IPW 和基于风险函数的 IPW 估计量都具有有界性，其次归因于基于风险函数的估计量比有界 IPW 估计量更有效。TMLE 的结果进一步支持这一推测：该估计量也不敏感于所采用的 PS 估计方法（图 16.2）。由于 TMLE 是代入估计量，它也像有界 IPW 与基于风险函数的 IPW 估计量一样，受益于有界性。基于风险函数的 IPW 估计量比有界 IPW 估计量具有更好的理论有效性，从而进一步降低了对 PS 估计误差的敏感性；同样，TMLE 比有界 IPW 估计量具有更好的有效性，也可解释为何它与基于风险函数的估计量一样，对 PS 估计误差不敏感。

表 16.2 比较了实践中采用各种估计选择时可得的推断；这些选择给出的结果均与先前随机试验一致。该表展示了理论上预期的有效性增益：基于风险函数的 IPW 估计量优于有界和无界 IPW 估计量。它还表明，随着 PS 估计方法越来越非参数化，基于风险函数的 IPW 估计量的估计标准误通常会减小。这种表面上的估计有效性增益，可由以下现象解释：PS 估计方法逐渐变得更灵活时，IPW 权重更加集中，大权重所占比例降低（结果见 Neugebauer 等，2016）。实际上，通过 SL 得到的稳定化逆概率权重竟然没有一个超过 30。这些结果与一种常见观点相矛盾。该观点认为，实践中更应采用基于模型的 PS 估计，因为数据自适应 PS 估计会揭示正值性假设的实际违反，产生更大的权重，从而增加 IPW 估计的变异性；而依据任意指定的参数模型估计逆概率权重，则可视为一种隐式权重截断方案，通过设定错误模型的平滑作用限制大权重比例。相反，本研究结果提示：当正值性假设实际上未被违反时，基于模型估计权重反而可能造成人为的正值性实际违反，即模型设定错误产生大权重。这种基于模型估计的实际后果，不仅是有偏推断，还包括更大的不确定性，而两者都可通过数据自适应 PS 估计避免。

实践中选择 TMLE 而不选择有界、无界或基于风险函数的 IPW 估计量，一个理由是 TMLE 的有效性可能带来估计精度增益。表 16.2 展示了本研究中的这种增益，但也表明，相比基于风险函数的 IPW 估计量，TMLE 的精度提升较小：在以 SL 估计行动机制时，3 年时六个累积风险差的基于风险函数 IPW 估计标准误，与 TMLE 估计标准误之比仅在 2\%—5\% 之间。由定义 TMLE 估计量影响曲线的式 (16.7) 可见，精度增益来自最小化干扰参数 $Q_0^\theta$ 估计的预测误差。因此，在协变量能很好预测结局的问题中，TMLE 的精度增益预期最大。表 16.2 报告的 TMLE 有效性增益较小，可能是由于测得的协变量中没有良好的结局预测变量，也可能是未充分利用这些协变量估计 $Q_0^\theta$，例如作出马尔可夫假设等模型设定错误。对后一种解释的考虑，促使我们使用 SL 灵活估计 $Q_0^\theta$，不依赖马尔可夫假设和任意的学习器选择，如 DSA 算法所考虑的仅主效应项 logistic 模型。表 16.1 展示了通过 SL 对干扰参数 $Q_0^\theta$ 作更非参数化估计所带来的有效性改善：当行动机制用 SL 估计时，3 年时六个累积风险差的基于风险函数 IPW 估计标准误，与 TMLE 估计标准误之比在 7\%—11\% 之间。

\begingroup
\scriptsize\setlength{\tabcolsep}{2.6pt}\renewcommand{\arraystretch}{1.08}
\begin{longtable}{@{}p{2.55cm}cccrrrrrr@{}}
\caption{六种 IPW 估计量与基于 DSA 的 TMLE 对 3 年时六个风险差（RD）的推断比较（12 个季度，即 $t_0=11$）。}\label{tab:ch16-two}\\
\toprule
RD 估计量 & $\theta_1$ & $\theta_2$ & $g_n^\theta$ & $\psi_n^{\theta_1,\theta_2}$ & $\Gamma_n^{\theta_1,\theta_2}$ & $\psi_n^{\theta_1,\theta_2,-}$ & $\psi_n^{\theta_1,\theta_2,+}$ & $p$ & RE\\
\midrule
\endfirsthead
\multicolumn{10}{l}{\small 表 16.2（续）}\\
\toprule
RD 估计量 & $\theta_1$ & $\theta_2$ & $g_n^\theta$ & $\psi_n^{\theta_1,\theta_2}$ & $\Gamma_n^{\theta_1,\theta_2}$ & $\psi_n^{\theta_1,\theta_2,-}$ & $\psi_n^{\theta_1,\theta_2,+}$ & $p$ & RE\\
\midrule
\endhead
\midrule\multicolumn{10}{r}{（续下页）}\\
\endfoot
\bottomrule
\endlastfoot
\shortstack[l]{Horvitz--Thompson\\$\psi^{\theta_1,\theta_2}_{n,\mathrm{HT}}$} & $8.5$ & $8$ & $g^\theta_{n,t,\mathrm{SL}}$ & $7.3\mathrm e{-03}$ & $4.3\mathrm e{-03}$ & $-1.1\mathrm e{-03}$ & $0.0156$ & $0.089$ & $1.085$\\
 & $8.5$ & $7.5$ & $g^\theta_{n,t,\mathrm{SL}}$ & $0.0231$ & $7.5\mathrm e{-03}$ & $8.4\mathrm e{-03}$ & $0.0378$ & $2\mathrm e{-03}$ & $1.078$\\
 & $8.5$ & $7$ & $g^\theta_{n,t,\mathrm{SL}}$ & $0.0506$ & $0.0129$ & $0.0254$ & $0.0758$ & $0$ & $1.082$\\
 & $8$ & $7.5$ & $g^\theta_{n,t,\mathrm{SL}}$ & $0.0158$ & $6.8\mathrm e{-03}$ & $2.5\mathrm e{-03}$ & $0.0292$ & $2\mathrm e{-02}$ & $1.078$\\
 & $8$ & $7$ & $g^\theta_{n,t,\mathrm{SL}}$ & $0.0433$ & $0.0128$ & $0.0182$ & $0.0685$ & $1\mathrm e{-03}$ & $1.083$\\
 & $7.5$ & $7$ & $g^\theta_{n,t,\mathrm{SL}}$ & $0.0275$ & $0.0118$ & $4.4\mathrm e{-03}$ & $0.0506$ & $0.019$ & $1.079$\\
\midrule
\shortstack[l]{有界 IPW\\$\psi^{\theta_1,\theta_2}_{n,\mathrm{bd}}$} & $8.5$ & $8$ & $g^\theta_{n,t,\mathrm{SL}}$ & $9.8\mathrm e{-03}$ & $4.6\mathrm e{-03}$ & $8\mathrm e{-04}$ & $0.0188$ & $0.032$ & $1.165$\\
 & $8.5$ & $7.5$ & $g^\theta_{n,t,\mathrm{SL}}$ & $0.0255$ & $7.9\mathrm e{-03}$ & $1\mathrm e{-02}$ & $0.041$ & $1\mathrm e{-03}$ & $1.138$\\
 & $8.5$ & $7$ & $g^\theta_{n,t,\mathrm{SL}}$ & $0.0572$ & $0.0127$ & $0.0322$ & $0.0822$ & $0$ & $1.072$\\
 & $8$ & $7.5$ & $g^\theta_{n,t,\mathrm{SL}}$ & $0.0157$ & $7.2\mathrm e{-03}$ & $1.5\mathrm e{-03}$ & $0.0298$ & $3\mathrm e{-02}$ & $1.144$\\
 & $8$ & $7$ & $g^\theta_{n,t,\mathrm{SL}}$ & $0.0474$ & $0.0127$ & $0.0224$ & $0.0723$ & $0$ & $1.072$\\
 & $7.5$ & $7$ & $g^\theta_{n,t,\mathrm{SL}}$ & $0.0317$ & $0.0117$ & $8.7\mathrm e{-03}$ & $0.0547$ & $7\mathrm e{-03}$ & $1.075$\\
\midrule
\shortstack[l]{基于风险函数的 IPW\\$\psi^{\theta_1,\theta_2}_{n,h}$} & $8.5$ & $8$ & $g_n^\theta$ & $0.0109$ & $5.7\mathrm e{-03}$ & $-3\mathrm e{-04}$ & $0.022$ & $0.056$ & \\
 & $8.5$ & $7.5$ & $g_n^\theta$ & $0.0212$ & $8.7\mathrm e{-03}$ & $4.1\mathrm e{-03}$ & $0.0382$ & $0.015$ & \\
 & $8.5$ & $7$ & $g_n^\theta$ & $0.0387$ & $0.0135$ & $0.0122$ & $0.0652$ & $4\mathrm e{-03}$ & \\
 & $8$ & $7.5$ & $g_n^\theta$ & $0.0103$ & $7.8\mathrm e{-03}$ & $-5\mathrm e{-03}$ & $0.0256$ & $0.187$ & \\
 & $8$ & $7$ & $g_n^\theta$ & $0.0278$ & $0.0138$ & $8\mathrm e{-04}$ & $0.0549$ & $0.044$ & \\
 & $7.5$ & $7$ & $g_n^\theta$ & $0.0175$ & $0.0133$ & $-8.5\mathrm e{-03}$ & $0.0435$ & $0.187$ & \\
\midrule
\shortstack[l]{基于风险函数的 IPW\\$\psi^{\theta_1,\theta_2}_{n,h}$} & $8.5$ & $8$ & $g^\theta_{n,t}$ & $9.9\mathrm e{-03}$ & $5.5\mathrm e{-03}$ & $-9\mathrm e{-04}$ & $0.0207$ & $0.071$ & \\
 & $8.5$ & $7.5$ & $g^\theta_{n,t}$ & $0.025$ & $7.8\mathrm e{-03}$ & $9.6\mathrm e{-03}$ & $0.0404$ & $1\mathrm e{-03}$ & \\
 & $8.5$ & $7$ & $g^\theta_{n,t}$ & $0.0418$ & $0.0126$ & $0.0172$ & $0.0664$ & $1\mathrm e{-03}$ & \\
 & $8$ & $7.5$ & $g^\theta_{n,t}$ & $0.0151$ & $7.1\mathrm e{-03}$ & $1.2\mathrm e{-03}$ & $0.029$ & $0.033$ & \\
 & $8$ & $7$ & $g^\theta_{n,t}$ & $0.0319$ & $0.0129$ & $6.7\mathrm e{-03}$ & $0.0571$ & $0.013$ & \\
 & $7.5$ & $7$ & $g^\theta_{n,t}$ & $0.0168$ & $0.012$ & $-6.7\mathrm e{-03}$ & $0.0403$ & $0.162$ & \\
\midrule
\shortstack[l]{基于风险函数的 IPW\\$\psi^{\theta_1,\theta_2}_{n,h}$} & $8.5$ & $8$ & $g^\theta_{n,t,\times}$ & $0.0129$ & $5.5\mathrm e{-03}$ & $2.1\mathrm e{-03}$ & $0.0237$ & $2\mathrm e{-02}$ & \\
 & $8.5$ & $7.5$ & $g^\theta_{n,t,\times}$ & $0.0337$ & $7.9\mathrm e{-03}$ & $0.0183$ & $0.0491$ & $0$ & \\
 & $8.5$ & $7$ & $g^\theta_{n,t,\times}$ & $0.0483$ & $0.0129$ & $0.0229$ & $0.0737$ & $0$ & \\
 & $8$ & $7.5$ & $g^\theta_{n,t,\times}$ & $0.0209$ & $6.9\mathrm e{-03}$ & $7.3\mathrm e{-03}$ & $0.0345$ & $3\mathrm e{-03}$ & \\
 & $8$ & $7$ & $g^\theta_{n,t,\times}$ & $0.0354$ & $0.0129$ & $1\mathrm e{-02}$ & $0.0608$ & $6\mathrm e{-03}$ & \\
 & $7.5$ & $7$ & $g^\theta_{n,t,\times}$ & $0.0146$ & $0.0121$ & $-9.2\mathrm e{-03}$ & $0.0383$ & $0.23$ & \\
\midrule
\shortstack[l]{基于风险函数的 IPW\\$\psi^{\theta_1,\theta_2}_{n,h}$} & $8.5$ & $8$ & $g^\theta_{n,t,\mathrm{SL}}$ & $8\mathrm e{-03}$ & $4\mathrm e{-03}$ & $1\mathrm e{-04}$ & $0.0159$ & $0.046$ & $1.025$\\
 & $8.5$ & $7.5$ & $g^\theta_{n,t,\mathrm{SL}}$ & $0.0216$ & $7.1\mathrm e{-03}$ & $7.7\mathrm e{-03}$ & $0.0356$ & $2\mathrm e{-03}$ & $1.024$\\
 & $8.5$ & $7$ & $g^\theta_{n,t,\mathrm{SL}}$ & $0.041$ & $0.0125$ & $0.0166$ & $0.0655$ & $1\mathrm e{-03}$ & $1.049$\\
 & $8$ & $7.5$ & $g^\theta_{n,t,\mathrm{SL}}$ & $0.0136$ & $6.5\mathrm e{-03}$ & $1\mathrm e{-03}$ & $0.0263$ & $0.035$ & $1.022$\\
 & $8$ & $7$ & $g^\theta_{n,t,\mathrm{SL}}$ & $0.033$ & $0.0124$ & $8.6\mathrm e{-03}$ & $0.0574$ & $8\mathrm e{-03}$ & $1.048$\\
 & $7.5$ & $7$ & $g^\theta_{n,t,\mathrm{SL}}$ & $0.0194$ & $0.0114$ & $-3\mathrm e{-03}$ & $0.0418$ & $9\mathrm e{-02}$ & $1.048$\\
\midrule
\shortstack[l]{基于 DSA 的 TMLE\\$\psi^{\theta_1,\theta_2,*}_{n}$} & $8.5$ & $8$ & $g^\theta_{n,t,\mathrm{SL}}$ & $9.1\mathrm e{-03}$ & $3.9\mathrm e{-03}$ & $1.4\mathrm e{-03}$ & $0.0168$ & $0.021$ & $1$\\
 & $8.5$ & $7.5$ & $g^\theta_{n,t,\mathrm{SL}}$ & $0.0238$ & $6.9\mathrm e{-03}$ & $0.0102$ & $0.0374$ & $1\mathrm e{-03}$ & $1$\\
 & $8.5$ & $7$ & $g^\theta_{n,t,\mathrm{SL}}$ & $0.0427$ & $0.0119$ & $0.0194$ & $0.0661$ & $0$ & $1$\\
 & $8$ & $7.5$ & $g^\theta_{n,t,\mathrm{SL}}$ & $0.0147$ & $6.3\mathrm e{-03}$ & $2.3\mathrm e{-03}$ & $0.0271$ & $2\mathrm e{-02}$ & $1$\\
 & $8$ & $7$ & $g^\theta_{n,t,\mathrm{SL}}$ & $0.0336$ & $0.0119$ & $0.0104$ & $0.0569$ & $5\mathrm e{-03}$ & $1$\\
 & $7.5$ & $7$ & $g^\theta_{n,t,\mathrm{SL}}$ & $0.0189$ & $0.0109$ & $-2.5\mathrm e{-03}$ & $0.0403$ & $0.083$ & $1$\\
\end{longtable}
\endgroup
\noindent{\footnotesize 对每个 RD，$g_0^\theta$ 的估计量、点估计、标准误估计、该标准误与基于 DSA 的 TMLE 标准误之比、95\% 置信区间上下限估计及相应 $p$ 值，分别记为 $g_n^\theta$、$\psi_n^{\theta_1,\theta_2}$、$\Gamma_n^{\theta_1,\theta_2}$、RE、$\psi_n^{\theta_1,\theta_2,-}$、$\psi_n^{\theta_1,\theta_2,+}$ 和 $p$。}\par


\begin{figure}[p]
\centering
\resizebox{.98\textwidth}{!}{\begin{tikzpicture}[x=1pt,y=-1pt]
\node[anchor=north west,inner sep=0] at (0,0) {\includegraphics[width=333.7pt]{figures/ch16_survival_base.pdf}};
\foreach \xx in {5.6,180.8}{
 \foreach \yy in {87.8,309.6}{
  \node[rotate=90,font=\fontsize{6.2}{7}\selectfont,inner sep=0] at (\xx,\yy) {生存概率 $P(T_d>t)$（权重未截断）};
 }
}
\foreach \xx in {103.4,277.6}{
 \foreach \yy in {191.9,414.2}{
  \node[font=\fontsize{7}{8}\selectfont,inner sep=0] at (\xx,\yy) {入组后的季度数};
 }
}
\end{tikzpicture}
}
\caption{每个面板展示对应于四种 TI 启动策略 $d_\theta$（$\theta=7,7.5,8,8.5$）的四条反事实生存曲线，在 16 个季度内的 TMLE 估计。左上、右上、左下、右下面板分别基于干扰参数 $g^\theta$ 的估计量 $g_n^\theta$、$g^\theta_{n,t}$、$g^\theta_{n,t,\times}$ 和 $g^\theta_{n,t,\mathrm{SL}}$ 得到。}
\label{fig:ch16-survival}
\end{figure}

再次强调，只有行动机制得到一致估计时，基于式 (16.7) 所定义影响曲线的 TMLE 推断才有有效性保证。为评价基于 IC 的 TMLE 推断表现，我们利用 10,000 个自助样本，评价图 16.2 中每个基于 DSA 的 TMLE 的变异性。为公平比较基于 IC 和基于自助法的 TMLE 变异性估计，我们在每个自助样本上，使用相同的已估计干扰参数构造 TMLE 点估计，即不在每个自助样本上重新计算 PS 估计量。这是因为，基于 IC 的 TMLE 估计量，仅在干扰参数 $g_0^\theta$ 已知时，才一致地估计 TMLE 标准误；而在行动机制得到一致估计时，它是保守的。当 TMLE 基于估计量 $g_n^\theta$、$g^\theta_{n,t}$ 或 $g^\theta_{n,t,\times}$ 时，3 年时六个累积风险差的基于 IC 的 TMLE 标准误估计，与自助法估计之比，显示基于 IC 的标准误估计存在明显高估，最高达 61\%（结果见 Neugebauer 等，2014a）。随着行动机制估计方法越来越非参数化，这些比值明显趋向理想值 1；使用 SL 估计行动机制时，全部比值基本都为 1。类似地，结果表明：对于使用较强参数形式估计量 $g_n^\theta$ 的有界 IPW 估计量，如果 IC 公式忽略归一化常数，自助法方差估计与基于 IC 的方差估计并不匹配；若在计算基于 IC 的估计时使用归一化常数，两者则相匹配。这些观察进一步支持以下判断：本研究中基于模型估计行动机制，产生了有偏的 PS 估计；而 SL 可用来防范这种参数模型设定错误。这些结果还表明，与前文对 TMLE 估计量本身的观察不同，基于 IC 的 TMLE 方差估计量对 PS 估计误差高度敏感。然而，在基于风险函数的 IPW 估计量的 IC 方差估计中，没有观察到这种敏感性（结果见 Neugebauer 等，2014a）。

\section{讨论}

我们展示了靶向学习在基于大型 EHR 数据库、具有时变暴露的真实世界比较效果研究中的可行性。我们说明了该方法如何成功调整时变混杂和选择偏差，并且相较于 IPW 估计获得有效性增益。对于 TMLE 和 IPW 两类估计量，我们均表明：相较于基于模型估计干扰参数，SL 估计可以显著减少偏差并提高有效性，而这种估计表现的改善以计算时间增加为代价。特别是，本章结果与“由于担心揭示正值性假设违反，实践中可能不宜采用 SL 估计”的观点相矛盾；相反，它们提示，实际常用的、基于任意参数模型的行动机制估计方法，在正值性假设事实上没有被违反时，反而可能造成人为违反，表现为极端的逆概率权重。

这些结果还提示，为使 TMLE 获得更显著的有效性增益，估计干扰参数 $Q_0^\theta$ 通常应依赖更加丰富的处理历史和协变量历史。原因在于，对行动机制某些组成部分可能成立的简化假设，如马尔可夫假设，通常不能预期推广到干扰参数 $Q_0^\theta$。若仍作出这些简化假设，就可能因 $Q_0^\theta$ 估计不一致，而使 TMLE 只能获得次优的有效性增益。

此外，本研究指出，有界性对于降低 TMLE 估计量对行动机制估计中轻微偏差的敏感性十分重要；但同时也提示，行动机制的数据自适应估计，对于获得适当的基于 IC 的 TMLE 推断至关重要，因为基于 IC 的 TMLE 标准误估计量，对行动机制估计中的轻微偏差高度敏感。这后一观察也提示，可使用自助法评价行动机制估计量的偏差，并指导选择实践中用于 TMLE 推断的 PS 估计量。例如，可以依次提高 $g_0^\theta$ 估计方法的数据自适应程度，直到基于 IC 与基于自助法的 TMLE 方差估计相匹配。另一种靶向性较弱的方法，是逐渐提高行动机制 SL 估计方法的数据自适应程度，直到交叉验证似然稳定；上述 PS 估计量选择准则可替代这一方法。

最后，当使用数据自适应方法估计干扰参数时，IPW 和 TMLE 估计量可能不再相对于本章用来推导保守方差估计量的那些影响曲线具有渐近线性。因此，当干扰参数采用 SL 估计，而非基于模型的极大似然估计时，实践中应用这些方差估计量可能缺少理论依据。近期，为在采用 SL 估计干扰参数时构造有效推断而开展的研究（van der Laan，2014b），促成了新的 IPW 与 TMLE 估计量（Benkeser 等，2017a）；它们具有渐近线性，且影响曲线已知。目前，该方法仅适用于单时点处理研究。还需要进一步研究，将这些 IPW 和 TMLE 方法推广到具有时变暴露和协变量的研究，如本章的强化治疗研究。在这一推广实现之前，有人可能主张：真实世界中具有时变干预的比较效果研究，其 TMLE 或 IPW 估计应依赖于基于模型的干扰参数估计，因为这样即可利用本章的保守方差估计量，在实践中得到理论上有效的推断。我们的反驳是：这些方差估计量只有在依据正确模型估计行动机制时才有效；若模型错误，所得置信区间的渐近覆盖率为 0。

\paragraph{致谢。}
作者感谢 HMO Research Network 的以下研究者，感谢他们提供各自研究点的数据，供本章二次分析使用：Denise M. Boudreau（Group Health）、Connie Trinacty（Kaiser Permanente Hawaii）、Gregory A. Nichols（Kaiser Permanente Northwest）、Marsha A. Raebel（Kaiser Permanente Colorado）以及 Kristi Reynolds（Kaiser Permanente Southern California）。本项目由美国卫生与公众服务部下属医疗研究与质量局依据合同 HHSA29020050016I 资助，是“建立证据以支持效果决策”（Developing Evidence to Inform Decisions about Effectiveness，DEcIDE）计划的一部分。本报告内容由作者负责。报告中的陈述不应被解释为获得医疗研究与质量局或美国卫生与公众服务部的认可。
\FloatBarrier
